\documentclass[12pt,a4paper]{article}

% Drake_preamble.tex - LaTeX Preamble Template
% 作者: 謝慕揚, MD, PhD, FESC
% 
% 使用說明:
% 1. 表格請使用 longtable，不要有垂直分隔線 |
% 2. 表格內換行使用 \newline，不要使用 \\
% 3. 藥物名稱、檢驗、檢查請維持英文
% 4. Reference 格式請使用 \href 加上驗證過的 DOI

% Including essential packages for Chinese typesetting
\usepackage{xeCJK}
\usepackage{fontspec}
% Configuring Chinese font with Noto Serif CJK TC, fallback to SimSun
\IfFontExistsTF{Noto Serif CJK TC}{
  \setCJKmainfont{Noto Serif CJK TC}
  \setCJKsansfont{Noto Serif CJK TC}
  \setCJKmonofont{Noto Serif CJK TC}
}{\setCJKmainfont{SimSun}
  \setCJKsansfont{SimSun}
  \setCJKmonofont{SimSun}
}

% Including other necessary packages
\usepackage{geometry}
\usepackage{titlesec}
\usepackage{enumitem}
\usepackage{xcolor}
\usepackage{colortbl}
\usepackage{tcolorbox}
\usepackage{booktabs}
\usepackage{longtable}
\usepackage{array}
\usepackage{graphicx}
\usepackage{tikz}
\usepackage{fancyhdr}
\usepackage{multicol}
\usepackage{datetime2}
\usepackage{hyperref}
\usepackage{mdframed}
\usepackage{amsmath}
\usepackage{amssymb}
\usepackage{multirow}

\hypersetup{
    colorlinks=true,
    linkcolor=emergency_blue,
    citecolor=emergency_blue,
    urlcolor=emergency_blue,
    pdfborder={0 0 0}
}

% Defining page geometry
\geometry{
  top=2.5cm,
  bottom=2.5cm,
  left=2cm,
  right=2cm,
  headheight=25pt
}

% Adjusting topmargin to compensate for increased headheight
\addtolength{\topmargin}{-10pt}

% Defining colors
\definecolor{emergency_red}{RGB}{186,24,27}
\definecolor{emergency_blue}{RGB}{0,114,188}
\definecolor{emergency_gray}{RGB}{120,120,120}
\definecolor{highlight_yellow}{RGB}{255,250,205}

% Setting title formats
\titleformat{\section}[block]{\Large\bfseries\color{emergency_red}}{\thesection}{10pt}{}
\titleformat{\subsection}[block]{\large\bfseries\color{emergency_blue}}{\thesubsection}{8pt}{}
\titleformat{\subsubsection}[block]{\normalsize\bfseries\color{emergency_gray}}{\thesubsubsection}{6pt}{}

% Defining custom environment for important notes
\newmdenv[
    linecolor=emergency_red,
    backgroundcolor=highlight_yellow,
    linewidth=2pt,
    topline=true,
    bottomline=true,
    leftline=true,
    rightline=true,
    innertopmargin=10pt,
    innerbottommargin=10pt,
    innerleftmargin=10pt,
    innerrightmargin=10pt
]{importantbox}

% Defining note command with tcolorbox
\newcommand{\note}[1]{
    \par
    \noindent
    \begin{tcolorbox}[
        colback=emergency_blue!5!white,
        colframe=emergency_blue,
        title=\textbf{重點提醒},
        fonttitle=\bfseries,
        boxrule=0.5pt,
        arc=3pt,
        left=6pt,
        right=6pt,
        top=6pt,
        bottom=6pt
    ]
    #1
    \end{tcolorbox}
}

% Key findings box
\newcommand{\keyfinding}[1]{
    \par
    \noindent
    \begin{tcolorbox}[
        colback=yellow!10!white,
        colframe=orange!75!black,
        title=\textbf{核心發現摘要},
        fonttitle=\bfseries,
        boxrule=0.5pt,
        arc=3pt
    ]
    #1
    \end{tcolorbox}
}

% Defining custom date format for ISO 8601
\DTMnewdatestyle{iso}{
  \renewcommand{\DTMdisplaydate}[4]{\number##1-\number##2-\number##3}
  \renewcommand{\DTMDisplaydate}[4]{\number##1-\number##2-\number##3}
}
\DTMsetdatestyle{iso}
\newcommand{\mydate}{\DTMtoday}


% Custom header for this document
\pagestyle{fancy}
\fancyhf{}
\fancyhead[L]{\textcolor{emergency_red}{Research Grant 教學講義}}
\fancyhead[R]{\textcolor{emergency_blue}{Carotid Web \& RBC Thrombosis}}
\fancyfoot[C]{\textcolor{emergency_gray}{\thepage}}
\renewcommand{\headrulewidth}{1pt}
\renewcommand{\headrule}{\hbox to\headwidth{\color{emergency_red}\leaders\hrule height \headrulewidth\hfill}}

\begin{document}

%============================================================
% Title Page
%============================================================
\begin{center}
{\LARGE\bfseries\textcolor{emergency_red}{Australia NHMRC Ideas Grant}}\\[0.3cm]
{\Large\textcolor{emergency_blue}{Research Proposal (2026-2030)}}\\[0.5cm]
{\Large\textbf{Beyond White Clots: Decoding RBC Mechanosensing in Carotid Web Thrombosis}}\\[0.3cm]
{\large 超越白色血栓：解碼紅血球在頸動脈蹼血栓形成中的機械感應機制}\\[0.5cm]
{\normalsize 整理：謝慕揚 MD, PhD, FESC \quad 日期：\mydate}
\end{center}

\vspace{0.5cm}

%============================================================
\section{研究背景與臨床問題}
%============================================================

\subsection{傳統血栓形成觀念}

\begin{itemize}[leftmargin=2em]
    \item \textbf{動脈血栓 (Arterial thrombosis)}：傳統認為主要由 \textbf{platelet-rich ``white clots''} 驅動
    \item \textbf{靜脈血栓 (Venous thrombosis)}：主要由 \textbf{RBC-rich ``red clots''} 組成
    \item 目前動脈血栓預防策略以 \textbf{antiplatelet therapy} 為主
\end{itemize}

\subsection{Carotid Web (頸動脈蹼) 的臨床挑戰}

\begin{importantbox}
\textbf{Carotid Web (CaW) 定義：}
\begin{itemize}[leftmargin=2em]
    \item 頸動脈球部 (carotid bulb) 後壁的薄膜狀突起
    \item 屬於 fibromuscular dysplasia 的一種變異型
    \item 造成血流動力學異常，形成 \textbf{recirculation zone}
    \item 與年輕患者 \textbf{cryptogenic stroke} 高度相關
\end{itemize}
\end{importantbox}

\textbf{臨床困境：}
\begin{itemize}[leftmargin=2em]
    \item CaW 患者即使接受標準 antiplatelet therapy，\textbf{復發性中風率仍然偏高}
    \item 傳統 antiplatelet 藥物無法有效預防此類血栓
    \item 提示這類血栓形成機制可能與一般動脈血栓不同
\end{itemize}

%============================================================
\section{核心假說 (Central Hypothesis)}
%============================================================

\keyfinding{
\textbf{研究假說：}在 Carotid Web 等複雜血管幾何結構中，\textbf{陡峭的剪切力梯度 (steep shear gradients)} 會活化紅血球上的 mechanosensitive ion channel \textbf{PIEZO1}。

PIEZO1 活化後觸發膜蛋白 \textbf{Band 3} 和 \textbf{Glycophorin A (GPA)} 的群聚 (clustering)，將紅血球轉變為具有黏附性的 \textbf{pro-thrombotic cells}，主動參與血栓形成。
}

\subsection{機制示意圖}

\begin{center}
\textbf{Carotid Web} $\rightarrow$ \textbf{Steep Shear Gradients} $\rightarrow$ \textbf{PIEZO1 Activation}

$\downarrow$

\textbf{Band 3/GPA Clustering} $\rightarrow$ \textbf{Pro-thrombotic RBCs} $\rightarrow$ \textbf{RBC-rich ``Red Clot''}
\end{center}

\begin{center}
\fbox{\parbox{13cm}{
\centering
\textbf{機制流程：}\\[0.3cm]
\textcolor{emergency_blue}{頸動脈蹼 (CaW)} $\xrightarrow{\text{血流動力學異常}}$
\textcolor{orange}{陡峭剪切力梯度} $\xrightarrow{\text{機械刺激}}$
\textcolor{red}{PIEZO1 活化}\\[0.2cm]
$\downarrow$ \small{Ca$^{2+}$ influx}\\[0.2cm]
\textcolor{emergency_red}{Band 3/GPA 群聚} $\rightarrow$
\textcolor{emergency_red}{促血栓紅血球} $\rightarrow$
\textbf{紅色血栓形成}
}}
\end{center}

%============================================================
\section{關鍵分子機制}
%============================================================

\subsection{PIEZO1 Mechanosensitive Ion Channel}

\begin{longtable}{p{4cm}p{10cm}}
\toprule
\textbf{特性} & \textbf{說明} \\
\midrule
\endfirsthead
基因 & \textit{PIEZO1} (2010年發現) \\
\midrule
功能 & Mechanosensitive cation channel，感應機械力並轉換為電化學訊號 \\
\midrule
在 RBC 的角色 & 調節細胞體積、感應剪切力、維持細胞變形能力 \\
\midrule
活化刺激 & Membrane stretch, shear stress, compression \\
\midrule
下游效應 & Ca$^{2+}$ influx → 細胞骨架重組 → 膜蛋白重新分布 \\
\midrule
臨床相關性 & PIEZO1 gain-of-function 突變與 hereditary xerocytosis 相關 \\
\bottomrule
\end{longtable}

\subsection{Band 3 與 Glycophorin A (GPA)}

\begin{longtable}{p{3cm}p{5.5cm}p{5.5cm}}
\toprule
\textbf{特性} & \textbf{Band 3} & \textbf{Glycophorin A (GPA)} \\
\midrule
\endfirsthead
別名 & Anion exchanger 1 (AE1), SLC4A1 & CD235a \\
\midrule
功能 & Cl$^-$/HCO$_3^-$ exchanger\newline 連接細胞膜與細胞骨架 & 主要唾液酸醣蛋白\newline 決定 MN 血型 \\
\midrule
在血栓中的角色 & 與 platelet 及 fibrin(ogen) 交互作用 & 介導 RBC-RBC 及 RBC-platelet 黏附 \\
\midrule
Clustering 效應 & 增加黏附位點暴露\newline 促進 RBC 聚集 & 增強細胞間交互作用 \\
\bottomrule
\end{longtable}

\note{
\textbf{PIEZO1-Band 3/GPA Axis：}當 PIEZO1 被機械力活化後，引發 Ca$^{2+}$ 內流，促使 Band 3 和 GPA 在細胞膜上重新分布並群聚。這種群聚增加了 RBC 表面的黏附位點，使原本「旁觀者」角色的紅血球轉變為主動參與血栓形成的促凝細胞。
}

%============================================================
\section{研究計畫架構 (Five Specific Aims)}
%============================================================

\subsection{Aim 1: Patient-specific Vessel-on-a-Chip 平台}

\begin{itemize}[leftmargin=2em]
    \item \textbf{技術：}利用患者 CT angiography 數據重建個人化「vessel-on-a-chip」
    \item \textbf{目的：}複製生理血流動力學條件
    \item \textbf{應用：}評估 CaW 患者的血栓形成風險
    \item \textbf{創新：}結合 3D 列印與微流體技術
\end{itemize}

\subsection{Aim 2 \& 3: 生物力學機制研究}

\begin{longtable}{p{3cm}p{11cm}}
\toprule
\textbf{方法} & \textbf{說明} \\
\midrule
\endfirsthead
Microfluidics & 模擬不同剪切力條件下的 RBC 行為 \\
\midrule
Biomembrane Force Probe (BFP) & 定量測量單一 RBC 與 platelet 之間的黏附力 \\
\midrule
PIEZO1-knockout mice & 驗證 PIEZO1 在血栓形成中的必要性 \\
\midrule
研究問題 & 機械應力與壓縮力如何透過 PIEZO1-Band 3/GPA 軸驅動 RBC-platelet 交互作用？ \\
\bottomrule
\end{longtable}

\subsection{Aim 4: 新型抑制劑開發}

\begin{importantbox}
\textbf{藥物開發目標：}
\begin{itemize}[leftmargin=2em]
    \item 開發針對 \textbf{PIEZO1} 的選擇性抑制劑
    \item 開發針對 \textbf{Band 3} 黏附功能的阻斷劑
    \item \textbf{關鍵要求：}阻斷促血栓交互作用，但\textbf{不影響正常止血功能}
\end{itemize}
\end{importantbox}

\subsection{Aim 5: 個人化藥物篩選平台}

\begin{itemize}[leftmargin=2em]
    \item 使用 patient-specific chip 平台測試新型化合物
    \item 與標準抗血栓藥物進行比較
    \item 使用臨床樣本進行驗證
    \item \textbf{目標：}預測個別患者的治療效果，實現精準醫療
\end{itemize}

%============================================================
\section{研究創新與臨床意義}
%============================================================

\subsection{突破傳統觀念}

\begin{center}
\begin{tabular}{p{6cm}p{7cm}}
\toprule
\textbf{傳統觀念} & \textbf{本研究提出} \\
\midrule
動脈血栓 = Platelet-rich white clot & 特定條件下形成 RBC-rich red clot \\
\midrule
RBC 在血栓中是被動旁觀者 & RBC 可主動參與血栓形成 \\
\midrule
Antiplatelet therapy 足以預防動脈血栓 & 需要 RBC-targeted therapy \\
\midrule
血栓預防策略一體適用 & 需要個人化風險評估 \\
\bottomrule
\end{tabular}
\end{center}

\subsection{預期臨床影響}

\begin{enumerate}[leftmargin=2em]
    \item \textbf{新治療標的：}PIEZO1 和 Band 3 成為抗血栓治療的新標的
    \item \textbf{更安全的抗凝策略：}針對 RBC 的治療可能比傳統抗凝血劑更安全
    \item \textbf{個人化醫療：}Vessel-on-a-chip 平台可預測個別患者的血栓風險與藥物反應
    \item \textbf{降低復發率：}為 CaW 患者提供更有效的中風預防策略
\end{enumerate}

%============================================================
\section{Carotid Web 臨床補充}
%============================================================

\subsection{流行病學}

\begin{itemize}[leftmargin=2em]
    \item 盛行率：一般人群約 1-2\%
    \item 在年輕 cryptogenic stroke 患者中可高達 \textbf{9-37\%}
    \item 好發於：年輕女性、非裔美國人
    \item 平均中風年齡：約 46-52 歲
\end{itemize}

\subsection{影像診斷}

\begin{longtable}{p{4cm}p{10cm}}
\toprule
\textbf{檢查方式} & \textbf{特徵} \\
\midrule
\endfirsthead
CTA (CT Angiography) & Axial view：carotid bulb 後壁薄膜狀突起\newline Sagittal view：「shelf-like」或「septum-like」結構 \\
\midrule
DSA (Digital Subtraction Angiography) & Lateral view：典型「web」外觀\newline 可見造影劑滯留 (contrast stasis) \\
\midrule
Duplex Ultrasound & 可見內膜不規則，但敏感度較低 \\
\bottomrule
\end{longtable}

\subsection{目前治療選項}

\begin{longtable}{p{4cm}p{5cm}p{5cm}}
\toprule
\textbf{治療} & \textbf{優點} & \textbf{缺點} \\
\midrule
\endfirsthead
Antiplatelet therapy & 非侵入性\newline 容易使用 & 復發率高 (約 30\%)\newline 可能無法針對 RBC-rich clot \\
\midrule
Anticoagulation & 可能對 red clot 更有效 & 出血風險\newline 缺乏 RCT 證據 \\
\midrule
Carotid stenting & 消除異常血流動力學\newline 預防復發 & 侵入性\newline 需長期 DAPT \\
\midrule
Carotid endarterectomy & 根本解決結構異常 & 手術風險\newline 可能過度治療 \\
\bottomrule
\end{longtable}

%============================================================
\section{Key Points 重點整理}
%============================================================

\begin{importantbox}
\begin{enumerate}[leftmargin=2em]
    \item \textbf{Carotid Web (CaW)} 是年輕患者 cryptogenic stroke 的重要病因，標準 antiplatelet therapy 預防效果不佳

    \item \textbf{核心假說：}CaW 的複雜血流動力學產生陡峭剪切力梯度，活化 RBC 上的 \textbf{PIEZO1} 機械感應通道

    \item \textbf{分子機制：}PIEZO1 活化 → Band 3/GPA clustering → RBC 轉變為 pro-thrombotic cell → RBC-rich ``red clot'' 形成

    \item \textbf{研究目標：}
    \begin{itemize}
        \item 建立 patient-specific vessel-on-a-chip 平台
        \item 驗證 PIEZO1-Band 3/GPA axis 機制
        \item 開發 RBC-targeted 抗血栓藥物
        \item 實現個人化血栓風險評估與治療
    \end{itemize}

    \item \textbf{臨床意義：}挑戰「動脈血栓 = 白色血栓」的傳統觀念，為 CaW 等高風險患者開發更有效、更安全的中風預防策略
\end{enumerate}
\end{importantbox}

%============================================================
\section{延伸閱讀}
%============================================================

\begin{enumerate}
    \item Choi PM, Singh D, Trivedi A, et al. Carotid Webs and Recurrent Ischemic Strokes in the Era of CT Angiography. \textit{AJNR Am J Neuroradiol}. 2015;36(11):2134-9.

    \item Coutinho JM, Derkatch S, Potber AJ, et al. Carotid artery web and ischemic stroke: A case-control study. \textit{Neurology}. 2017;88(1):65-69.

    \item Coste B, Mathur J, Schmidt M, et al. Piezo1 and Piezo2 are essential components of distinct mechanically activated cation channels. \textit{Science}. 2010;330(6000):55-60.

    \item Cahalan SM, Lukacs V, Ranade SS, et al. Piezo1 links mechanical forces to red blood cell volume. \textit{eLife}. 2015;4:e07370.

    \item Byrnes JR, Bhatt A, Daly ME, et al. Red blood cells in thrombosis. \textit{Blood}. 2018;130(16):1795-1799.
\end{enumerate}

\end{document}
